This thesis presents a stateful Fault Detection, Isolation, and Recovery (FDIR) framework for Service-Oriented Architecture (SOA) based industrial routers. Existing monitoring and recovery approaches for distributed systems treat fault detection, dependency analysis, and recovery as separate concerns and do not maintain runtime state that distinguishes new fault occurrences from ongoing violations. This results in redundant alerts, repeated recovery actions, and uncoordinated handling of cascading faults — conditions that are incompatible with the determinism and resource constraints of embedded industrial platforms.

The proposed framework addresses four gaps identified in the literature: (1) the absence of runtime fault episode tracking, (2) the lack of dependency-aware recovery coordination, (3) the missing enforcement of at-most-once recovery execution per fault episode, and (4) the reliance on computational methods unsuitable for embedded platforms. The framework integrates structured telemetry ingestion, a directed service dependency graph, rule-based fault detection with rising-edge transition tracking, a fault episode lifecycle model, root-cause-only resolution with transitive suppression of downstream symptoms, and persistent fault escalation into a unified seven-phase processing pipeline.

A prototype implementation (`prism-police`) on the Belden/Hirschmann PRISM industrial router platform validates the design through 47 unit tests covering all eight functional and nine of eleven non-functional requirements. Quantitative single-run experiments comparing the prototype with Prometheus and Nagios demonstrate that the episode model reduces alert volume by two to five orders of magnitude under sustained fault conditions, while CPU utilisation remains at 77.49% of a single core versus 313.83% (Prometheus) and 1,035.97% (Nagios) at 1,000 samples per second. Memory consumption stays between 13.92 and 29.37 MB across all tested scenarios. The evaluation is limited to a single platform and desktop-class hardware; production-platform validation and long-duration resource profiling remain open.